We test what a candidate-centered band of native CT can add to lung nodule ranking. At each released proposal we preserve a $10\mm$ ball and restore a band around it before interpolation. On \EvalScans{} reused LUNA16 scans, the primary two-scale 3D ranker rises from \BallCpm\% to \BandCpm\% CPM, a gain of \RestorePrimary{} points with a multiplicity-adjusted patient interval. A second fit gains \RestoreRepeat{} points with a pointwise interval. Matched sensitivity at 1 false positive per scan rises by \PrimaryOneGain{} and \RepeatOneGain{} points. Two controls clarify the dense advantage. Relative angular attention on the same 384 exterior readings leaves a multi-point gap. Even when the exact native sources of those readings enter the same image encoder as the band, the band leads by \CoverageGap{} CPM points. The matched sensitivity gains are largest among lesions whose majority-reader contour can be reached by the band at a matching proposal. A distinct single-stream 3D model also gains \SingleEstimate{} CPM points. These results locate useful evidence in the added native coverage and show that a proposal-centered exterior can restore target tissue as well as surrounding anatomy.
